\documentclass{article}

% -------------------------------------------------
% Packages
% -------------------------------------------------

\usepackage{xcolor}
\usepackage{amsmath}
\usepackage{amssymb}
\usepackage{amsthm}
\usepackage{graphicx}
\usepackage{subcaption}
\usepackage{matlab-prettifier}
\usepackage{listings}
\usepackage{blindtext}
\usepackage{geometry}
\usepackage{float}
\usepackage{dirtytalk}
\usepackage{enumitem}
\usepackage{derivative}
\usepackage{biblatex}
\usepackage{svg}
\usepackage{hyperref}
\usepackage{bbm}

% -------------------------------------------------
% Colors
% -------------------------------------------------

\definecolor{darkgreen}{rgb}{0.0, 0.5, 0.0}

% -------------------------------------------------
% Hyperlinks
% -------------------------------------------------

\hypersetup{
    colorlinks=true,
    linkcolor=blue,
    citecolor=blue,
    urlcolor=blue
}

% -------------------------------------------------
% Theorem environments
% -------------------------------------------------

% Numbered
\newtheorem{theorem}{Theorem}[section]
\newtheorem{lemma}[theorem]{Lemma}
\newtheorem{proposition}[theorem]{Proposition}

% Unnumbered
\newtheorem*{theorem*}{Theorem}
\newtheorem*{lemma*}{Lemma}
\newtheorem*{proposition*}{Proposition}

\theoremstyle{definition}
\newtheorem{definition}[theorem]{Definition}
\newtheorem{example}[theorem]{Example}

\theoremstyle{remark}
\newtheorem{remark}[theorem]{Remark}

% -------------------------------------------------
% Equation numbering
% -------------------------------------------------

% Reset equation numbering at each subsection
\makeatletter
\@addtoreset{equation}{subsection}
\makeatother

% If you want equations to be numbered by subsection,
% uncomment the following line:
% \renewcommand{\theequation}{\thesubsection.\arabic{equation}}

% -------------------------------------------------
% Python / MATLAB code formatting
% -------------------------------------------------

\lstset{
    language=Python,
    basicstyle=\ttfamily\small,
    keywordstyle=\color{blue},
    stringstyle=\color{darkgreen},
    commentstyle=\color{gray},
    showstringspaces=false,
    frame=single,
    numbers=left,
    numberstyle=\tiny\color{gray},
    breaklines=true
}

% -------------------------------------------------
% Page geometry
% -------------------------------------------------

\geometry{
    a4paper,
    total={160mm,250mm},
    left=25mm,
    top=30mm,
    footskip=10mm
}

% -------------------------------------------------
% Document settings
% -------------------------------------------------

\setlength{\parindent}{0pt}

% -------------------------------------------------
% Title
% -------------------------------------------------

\title{Monte Carlo Simulation Uge 3}
\author{Thomas Sejer Canham}
\date{\today}

% =================================================
% Document
% =================================================

\begin{document}

% -------------------------------------------------
% Title page
% -------------------------------------------------

\begin{titlepage}

    \maketitle
    \thispagestyle{empty}

\end{titlepage}

\section{Sampling Uniformly from Regions}
This section generalizes the idea of sampling uniformly from a region in $\mathbb{R}^n$. The goal is to sample points uniformly from a given region $B \subseteq \mathbb{R}^n$. Suppose that we sample $X = (X_1, X_2, \ldots, X_n)$ uniformly from a region $B$. Then the joint density of $X$ is given by
\begin{equation}
    f_{X_1X_2\ldots X_n}(x_1, x_2, \ldots, x_n) = \dfrac{1}{|B|},
\end{equation}
where $|B|$ is the volume (or area, in the case of $\mathbb{R}^2$) of the region $B$. The deinsity is constant in the region $B$ and zero outside of $B$. We can use this to find marginal densities of the individual components $X_i$. For example, the marginal density of $X_1$ is given by
\begin{equation}
    f_{X_1}(x_1) = \int_{B_{x_1}} \int_{-\infty}^{\infty} \ldots \int_{-\infty}^{\infty} f_{X_1X_2\ldots X_n}(x_1, x_2, \ldots, x_n) \, dx_2 \ldots dx_n = \int_{B_{x_1}} \dfrac{1}{|B|} \, dx_2 \ldots dx_n = \dfrac{|B_{x_1}|}{|B|}.
\end{equation}
This means that the marginal density of $X_1$ is proportional to the volume of the region $B_{x_1}$, which is the slice of the region $B$ at a fixed value of $x_1$. This can be generalized to any component $X_i$.

\subsection{Task 12 Acceptance-Rejection}
We illustrate using Task 12 as an example. 
We are given the distribution function $f(x) = C  x^{-3/2} \dfrac{1}{\sqrt{2\pi}} e^{-(x-1)^2/2x}$, and we wish to find a dominating distribution $g(x)$ such that $f(x) \leq M g(x)$ for some constant $M$. If we expand the expression for $f(x)$ and disregard the constant, as this can always be absorbed into the constant $M$, we have that
\begin{equation}
    f(x) \propto x^{3/2} e^{-(x-1)^2/2x} = x^{-3/2} e^{-x/2} e^{1/2} e^{-1/2x} \propto x^{-3/2} e^{-x/2} e^{-1/2x}.
\end{equation}
The behaviour near $x= 0$ can be assessed graphically, however the tail behaviour as $x \to \infty$ can be assessed analytically. We have that $x^{-3/2} \to 0$ as $x \to \infty$, and $e^{-1/2x} \to 1$ as $x \to \infty$. So we have to match the tail behaviour of $e^{-x/2}
$. Thus if we choose the exponential distribution $g(x) = \lambda e^{-\lambda x}$, we have that $f(x) \leq M g(x)$ for some constant $M$ and $\lambda = 1/2$. Graphically, we can see that $M \approx 2.7$.\vspace{1em}



Let $(X,Y)$ be uniformly distributed on the region $B = \{(x,y) \in \mathbb{R}^2 : 0 \leq y \leq f(x)\}$. Where $f(x)$ is a distribution function. We then wish to find the marginal density of $X$. We have that
\begin{equation}
    |B| = \int \int_B \mathbbm{1}(0 \leq y \leq f(x)) \, dy \, dx = \int_{-\infty}^{\infty} \int_0^{f(x)} \, dy \, dx = \int_{-\infty}^{\infty} f(x) \, dx = 1,
\end{equation}
Now for a fixed $x$, we have that
\begin{equation}
    f_{X}(x) = \dfrac{\int_0^{f(x)} f_{XY}(x,y)\, dy}{|B|} = \dfrac{\int_0^{f(x)} 1\, dy}{1} = f(x).
\end{equation}
This shows us that if we sample uniformly from the region below the curve of a distribution function $f(x)$, then the marginal density of $X$ is given by $f(x)$. This is the basis for the acceptance-rejection method.

\section{Transformation of Variables}
Suppose that we have a vector of random variables $X = (X_1, X_2, \ldots, X_n)$ with joint density $f_X(x)$. We wish to find the joint density of a new vector of random variables $Y = (Y_1, Y_2, \ldots, Y_n)$, where $Y = g(X)$ for some function $g$. If $g$ is a one-to-one transformation, then we can find the joint density of $Y$ using the change of variables formula. Let $h = g^{-1}$ be the inverse transformation. Then the joint density of $Y$ is given by
\begin{equation}
    f_Y(y) = f_X(h(y)) \left| \det J_h(y) \right|,
    \quad y \in \mathbb{R}^n.
\end{equation}
\subsection{Task 13}
We are given the transformation of $(X,Y)$ to $(V,U)$, where $X = \sqrt{V}cos(U)$ and $Y = \sqrt{V}sin(U)$, Where $ V > 0$ and $0 < U < \pi/2$. $X$ and $Y$ independent and standard normal. We wish to find the joint density of $(V,U)$. Note that the trasform implies that
$$X^2 + Y^2 = V$$
And since $(X,Y)$ is jointly standard normal, we have that
$$f_{X,Y}(x,y) = \frac{1}{2\pi} e^{-\frac{1}{2}(x^2 + y^2)} = \frac{1}{2\pi} e^{-\frac{1}{2}V}$$
\begin{equation}
    h(v,u) = (\sqrt{v}cos(u), \sqrt{v}sin(u)) = (x,y),
\end{equation}
and the Jacobian matrix of $h$ is
\begin{equation}
    J_h(v,u) = \begin{pmatrix}
        \frac{\partial x}{\partial v} & \frac{\partial x}{\partial u} \\
        \frac{\partial y}{\partial v} & \frac{\partial y}{\partial u}
    \end{pmatrix} = \begin{pmatrix}
        \frac{1}{2\sqrt{v}}cos(u) & -\sqrt{v}sin(u) \\
        \frac{1}{2\sqrt{v}}sin(u) & \sqrt{v}cos(u)
    \end{pmatrix}.
\end{equation}
The determinant of the Jacobian matrix is
\begin{equation}
    \det J_h(v,u) = \frac{1}{2\sqrt{v}}cos(u) \cdot \sqrt{v}cos(u) - (-\sqrt{v}sin(u)) \cdot \frac{1}{2\sqrt{v}}sin(u) = \frac{1}{2}cos^2(u) + \frac{1}{2}sin^2(u) = \frac{1}{2}.
\end{equation}
Therefore, the joint density of $(V,U)$ is
\begin{equation}
    f_{V,U}(v,u) = f_{X,Y}(h(v,u)) \left| \det J_h(v,u) \right| =  \frac{1}{2\pi} e^{-\frac{1}{2}V} \cdot \frac{1}{2}.
\end{equation}
We see that this is a product of two functions, one in $v$ and one in $u$, which means that $V$ and $U$ are independent. More specifically, we have that $V$ is exponentially distributed with rate parameter $\lambda = 1/2$, and $U$ is uniformly distributed on the interval $(0, \pi/2)$. This then means that (V,U) are points at distance $V$ from the origin, and the angle $U$ is uniformly distributed on the interval $(0, \pi/2)$. \vspace{2em}

Suppose we want point uniformly from the unit sphere. Then we are interested in points $(A,B,C)$ such that $A^2 + B^2 + C^2 = 1$. We have just seen above that independant standard normals are direction invariant, since the joint density functions becomes a function of the random variables through $||(X,Y,Z)||^2$. As such if we normalize the vector $(X,Y,Z)$, such that $(X,Y,Z) \to \frac{(X,Y,Z)}{||(X,Y,Z)||}$, then we have a point uniformly distributed on the unit sphere, where $A = \frac{X}{||(X,Y,Z)||}$, $B = \frac{Y}{||(X,Y,Z)||}$, and $C = \frac{Z}{||(X,Y,Z)||}$. If we substitute into the equation for the unit sphere, we have that
\begin{equation}
    A^2 + B^2 + C^2 = \frac{X^2}{||(X,Y,Z)||^2} + \frac{Y^2}{||(X,Y,Z)||^2} + \frac{Z^2}{||(X,Y,Z)||^2} = \frac{X^2 + Y^2 + Z^2}{||(X,Y,Z)||^2} = 1.
\end{equation}

\section{Task 15}
In the following we are interested in finding a bounding box for the region $B = \{(x_1,x_2) \in \mathbb{R}^2 : 0 \leq x_1 \leq \sqrt{g(x_1/x_2)}\}$, where $f(x)$ is a distribution function. We have a given candidate density function $g(x) = xe^{-x}$. The steps we need to follow are:
\begin{enumerate}
    
    \item Find the supremum of $\sqrt{g(x)}$, this will give us the upper bound for $x_1$.
    \item  Define $x = x_2/x_1$, note that $x\in \mathbb{R}$, can be both negative and positive.
    \item Find the bounds for $x_2$, since $x_2 = xx_1$, it follows from the inequality in $B$, and the fact that $x$ can be both negative and positive, that $ -sup_{x\in \mathbb{R}} x \sqrt{g(x)}< x_2 < sup_{x\in \mathbb{R}} x \sqrt{g(x)}$.
\end{enumerate}
For the the function $g(x) = xe^{-x}$, we find the supremum of $\sqrt{g(x)}$ to be $\frac{1}{\sqrt{e}}$. For the bounds of $x_2$, we have that $x_2 = xx_1$, and since $x$ can be both negative and positive, we have that $ -sup_{x\in \mathbb{R}} x \sqrt{g(x)}< x_2 < sup_{x\in \mathbb{R}} x \sqrt{g(x)}$. We find the supremum of $x\sqrt{g(x)}$ to be $\approx 1.16$. Thus we have that the bounding box is given by
\begin{equation}
    A = \{(x_1,x_2) \in \mathbb{R}^2 : 0 \leq x_1 \leq \frac{1}{\sqrt{e}}, -1.16 < x_2 < 1.16\}.
\end{equation}
Thus we need to sample from points $U \sim Unif(0, \frac{1}{\sqrt{e}})$ and $V \sim Unif(-1.16, 1.16)$, and accept the points that fall in the region $B$. Note for this perticular function, negative values of $x_2$ will create point in the complex plane, so we can further restrict the bounding box to be $A = \{(x_1,x_2) \in \mathbb{R}^2 : 0 \leq x_1 \leq \frac{1}{\sqrt{e}}, 0 < x_2 < 1.16\}$.

\end{document}
